\documentclass[10pt,a4paper]{amsart}
\usepackage{amsmath,amssymb,amsthm,graphicx}
\usepackage{verbatim}
\newenvironment{reference}{\comment}{\endcomment}
%\newenvironment{reference}{}{}
\newenvironment{slogan}{\comment}{\endcomment}
\newenvironment{history}{\comment}{\endcomment}

% For commutative diagrams we use Xy-pic
\usepackage[all]{xy}

% We use 2cell for 2-commutative diagrams.
\xyoption{2cell}
% Theorem environments.
%
\theoremstyle{plain}
\newtheorem{theorem}[subsection]{Theorem}
\newtheorem{proposition}[subsection]{Proposition}
\newtheorem{lemma}[subsection]{Lemma}

\theoremstyle{definition}
\newtheorem{definition}[subsection]{Definition}
\newtheorem{example}[subsection]{Example}
\newtheorem{exercise}[subsection]{Exercise}
\newtheorem{situation}[subsection]{Situation}

\theoremstyle{remark}
\newtheorem{remark}[subsection]{Remark}
\newtheorem{remarks}[subsection]{Remarks}

\numberwithin{equation}{subsection}

% Macros
%
\def\lim{\mathop{\mathrm{lim}}\nolimits}
\def\colim{\mathop{\mathrm{colim}}\nolimits}
\def\Spec{\mathop{\mathrm{Spec}}}
\def\Hom{\mathop{\mathrm{Hom}}\nolimits}
\def\Ext{\mathop{\mathrm{Ext}}\nolimits}
\def\SheafHom{\mathop{\mathcal{H}\!\mathit{om}}\nolimits}
\def\SheafExt{\mathop{\mathcal{E}\!\mathit{xt}}\nolimits}
\def\Sch{\mathit{Sch}}
\def\Mor{\mathop{\mathrm{Mor}}\nolimits}
\def\Ob{\mathop{\mathrm{Ob}}\nolimits}
\def\Sh{\mathop{\mathit{Sh}}\nolimits}
\def\NL{\mathop{N\!L}\nolimits}
\def\CH{\mathop{\mathrm{CH}}\nolimits}
\def\proetale{{pro\text{-}\acute{e}tale}}
\def\etale{{\acute{e}tale}}
\def\QCoh{\mathit{QCoh}}
\def\Ker{\mathop{\mathrm{Ker}}}
\def\Im{\mathop{\mathrm{Im}}}
\def\Coker{\mathop{\mathrm{Coker}}}
\def\Coim{\mathop{\mathrm{Coim}}}

% Boxtimes
%
\DeclareMathSymbol{\boxtimes}{\mathbin}{AMSa}{"02}

%
% Macros for moduli stacks/spaces
%
\def\QCohstack{\mathcal{QC}\!\mathit{oh}}
\def\Cohstack{\mathcal{C}\!\mathit{oh}}
\def\Spacesstack{\mathcal{S}\!\mathit{paces}}
\def\Quotfunctor{\mathrm{Quot}}
\def\Hilbfunctor{\mathrm{Hilb}}
\def\Curvesstack{\mathcal{C}\!\mathit{urves}}
\def\Polarizedstack{\mathcal{P}\!\mathit{olarized}}
\def\Complexesstack{\mathcal{C}\!\mathit{omplexes}}
% \Pic is the operator that assigns to X its picard group, usage \Pic(X)
% \Picardstack_{X/B} denotes the Picard stack of X over B
% \Picardfunctor_{X/B} denotes the Picard functor of X over B
\def\Pic{\mathop{\mathrm{Pic}}\nolimits}
\def\Picardstack{\mathcal{P}\!\mathit{ic}}
\def\Picardfunctor{\mathrm{Pic}}
\def\Deformationcategory{\mathcal{D}\!\mathit{ef}}
\usepackage{cleveref}
\crefname{originaldefinition}{original definition}{original definitions}
\Crefname{originaldefinition}{Original definition}{Original definitions}
\makeatletter
\newlabel{definition-property-morphisms-local}{{03YH}{source}{}{sourcepointer-1}{}}
\newlabel{definition-property-morphisms-local@cref}{{[originaldefinition][0][]03YH}{[1][0][]source}{}{}{}}
\makeatother
\title{Criterion for a property of morphisms of algebraic spaces to be fpqc local on the base}
\author{The Stacks Project Authors}
\date{stacks-project:spaces-descent.tex:local-snapshot-20261012; GNU FDL 1.2 or later}
\begin{document}
\maketitle
\noindent\textit{Editorial note.} All delivered mathematical text is the author's own, copied verbatim from the pinned Stacks Project chapter file spaces-descent.tex; no statement or proof step was rewritten, reordered or removed.\par\medskip
\noindent\textit{Editorial note.} The counted span cites one Topologies on Spaces lemma plus same-chapter items, recorded in assembly.external\_references with canonical tags as navigation pointers.\par\medskip
\noindent\textit{Allowed standard prerequisites (editorial).}\begin{itemize}
\item Definition of a property fpqc local on the base. Stacks Project, Descent on Spaces, Definition definition-property-morphisms-local. Supplies the target notion.
\item Standard fpqc coverings refine fpqc coverings. Stacks Project, Topologies on Spaces, Lemma spaces-topologies-lemma-fpqc. Supplies the standard affine refinement \$\textbackslash{}\{Z\_j \textbackslash{}to Z\textbackslash{}\}\$ with maps into the \$Y\_\{i\_j\}\$.
\end{itemize}\medskip
\ifdefined\hypertarget\hypertarget{sourcepointer-1}{}\fi
\noindent\textit{Source reference.} Original definition 03YH: sources/stacks-chow/tags.txt, source lines 5184--5184. Cross-chapter standard prerequisite cited by the author inside the counted span; it is used as a pointer and its own proof is not part of this item.\par\medskip
\begin{lemma}
\label{lemma-descending-properties-morphisms}
Let $S$ be a scheme. Let $\mathcal{P}$ be a property of morphisms of
algebraic spaces over $S$. Assume
\begin{enumerate}
\item if $X_i \to Y_i$, $i = 1, 2$ have property $\mathcal{P}$ so
does $X_1 \amalg X_2 \to Y_1 \amalg Y_2$,
\item a morphism of algebraic spaces $f : X \to Y$ has property
$\mathcal{P}$ if and only if for every affine scheme $Z$ and
morphism $Z \to Y$ the base change $Z \times_Y X \to Z$ of $f$
has property $\mathcal{P}$, and
\item for any surjective flat morphism of affine schemes
$Z' \to Z$ over $S$ and a morphism $f : X \to Z$ from an algebraic space
to $Z$ we have
$$
f' : Z' \times_Z X \to Z'\text{ has }\mathcal{P}
\Rightarrow
f\text{ has }\mathcal{P}.
$$
\end{enumerate}
Then $\mathcal{P}$ is fpqc local on the base.
\end{lemma}

\begin{proof}
If $\mathcal{P}$ has property (2), then it is automatically
stable under any base change. Hence the direct implication in
Definition \ref{definition-property-morphisms-local}.

\medskip\noindent
Let $\{Y_i \to Y\}_{i \in I}$ be an fpqc covering of algebraic spaces over $S$.
Let $f : X \to Y$ be a morphism of algebraic spaces over $S$.
Assume each base change $f_i : Y_i \times_Y X \to Y_i$ has property
$\mathcal{P}$. Our goal is to show that $f$ has $\mathcal{P}$.
Let $Z$ be an affine scheme, and let $Z \to Y$ be a morphism.
By (2) it suffices to show that the morphism of algebraic spaces
$Z \times_Y X \to Z$ has $\mathcal{P}$.
Since $\{Y_i \to Y\}_{i \in I}$ is an fpqc covering we know there
exists a standard fpqc covering $\{Z_j \to Z\}_{j = 1, \ldots , n}$
and morphisms $Z_j \to Y_{i_j}$ over $Y$ for suitable indices $i_j \in I$.
Since $f_{i_j}$ has $\mathcal{P}$ we see that
$$
Z_j \times_Y X
=
Z_j \times_{Y_{i_j}} (Y_{i_j} \times_Y X)
\longrightarrow
Z_j
$$
has $\mathcal{P}$ as a base change of $f_{i_j}$ (see first remark of the
proof). Set $Z' = \coprod_{j = 1, \ldots, n} Z_j$, so that $Z' \to Z$ is
a flat and surjective morphism of affine schemes over $S$. By (1)
we conclude that $Z' \times_Y X \to Z'$ has property $\mathcal{P}$.
Since this is the base change of the morphism $Z \times_Y X \to Z$
by the morphism $Z' \to Z$ we conclude that
$Z \times_Y X \to Z$ has property $\mathcal{P}$ as desired.
\end{proof}

\end{document}
